\section{Controller Tests}\label{sec:tests}
This section reports what the built controller did on the system, measured against the four
specifications of Section~\ref{sec:design}.

\subsection{Tests performed}
Four tests cover the specifications, and each one produces a single number that decides a pass
or a fail.

\begin{enumerate}
    \item \textbf{Setpoint step.} The setpoint is stepped and the response is logged until it
    settles. The trace gives both the steady-state offset as a fraction of the step, which tests
    the $8$ percent tracking specification, and the peak against the final value, which tests
    the $30$ percent overshoot specification.
    \item \textbf{Repeated steps at different setpoints.} Confirms the behaviour is not specific
    to one operating point, and exposes any asymmetry between upward and downward steps caused
    by the input clipping at $0$ and $5\ \mathrm{V}$.
    \item \textbf{Gain variation.} The plant cannot be varied, so the trimpot is moved $\pm 10$
    percent about the design value and the step is repeated at each extreme. This covers the
    component tolerance specification directly and stands in for the $10$ percent uncertainty in
    the aerodynamic constant, since both enter the closed-loop denominator through the product
    $K_p k A$.
    \item \textbf{Demonstration.} The standard step testing process run with Mr De Maar, with a
    screenshot of the controller tracking the setpoint recorded for
    Appendix~\ref{sec:appendix}.
\end{enumerate}

Each test has a predicted value computed before the bench session, so that the comparison is a
test of the design rather than a description of whatever happened. Table~\ref{tab:predictions}
lists them.

\begin{table}[htbp]
\centering
\caption{Predicted results at $K_p = 0.005$ and a $7\ \mathrm{m}$ setpoint, from the simulation
of Section~\ref{sec:design}. A bench result outside these bands means the model is wrong, not the
controller.}
\label{tab:predictions}
\begin{tabular}{lrrr}
\toprule
Quantity & Nominal & At $K_p -10\%$ & At $K_p +10\%$ \\
\midrule
Steady-state offset (\%) & 6.0 & 7.5 & 5.0 \\
Overshoot (\%)           & 17.3 & 13.0 & 21.0 \\
$2$ percent settling time (s) & 76 & 61 & 71 \\
Peak height (m)          & 8.21 & 7.91 & 8.47 \\
Command range (V)        & \multicolumn{3}{c}{2.49 to 2.54, no saturation} \\
\bottomrule
\end{tabular}
\end{table}

The command range matters as much as the response. The design only holds while the command stays
inside $0$ to $5\ \mathrm{V}$, and at this gain it moves barely $50\ \mathrm{mV}$ from the hover
bias, so any clipping seen on the bench would indicate a wiring or bias fault rather than a
design limit.

% Awaiting session 2: logged step responses, the measured offset and overshoot at each gain
% setting, and the demo screenshot.

\subsection{Performance evaluation}
Each specification is reported as a pass or a fail with the measured number that decides it.

\begin{table}[htbp]
\centering
\begin{tabular}{llll}
\toprule
\textbf{Specification} & \textbf{Test} & \textbf{Measured} & \textbf{Verdict} \\
\midrule
Tracking error below 8\%     & Setpoint step       & & \\
Overshoot below 30\%         & Same trace          & & \\
10\% uncertainty in $b$      & Gain varied $\pm$10\% & & \\
10\% component tolerance     & Gain varied $\pm$10\% & & \\
\bottomrule
\end{tabular}
\caption{Measured performance against each specification. Completed from the session 2 traces.}
\label{tab:spec_results}
\end{table}

Any departure from the prediction is attributed to one of three causes: an error in the identified
parameters, a difference between the designed and the built gain, or input saturation taking the
loop outside the linear range. Saturation is checked first, since it is the only one that can be
ruled in or out from the command trace alone.
